\section{The four Sundays}
\label{sec:pc-advent-sundays}

Find the Sunday in your parish calendar. Its letter, A, B or C, stays the
same all Advent. If the calendar gives only the civil date, take the year
of the January about to begin. A year divisible by three gives C, the
next year A, and the next B. Advent begins on the Sunday nearest November
30. Count that Sunday as the first. A reader arriving on a weekday can go
straight to \hyperref[sec:pc-advent-weekdays]{the weekday readings}.

On each line the references follow the order heard at Mass. The first
reading leads to the responsorial psalm or canticle, then to the second
reading and the Gospel. The numbers in parentheses identify the Mass in
the U.S. \emph{Lectionary for Mass}. Psalm numbers in the lists follow
the U.S. Lectionary. Other Catholic Bibles sometimes divide chapters
and verses differently as well. The short prayers for home are from the
Douay--Rheims Bible. Its psalm numbers and several places where its
prophetic readings fall differently are marked here.

\subsection{First Sunday. Who is coming?}

\begin{description}
\item[Year A (1)] Isaiah 2:1--5. Psalm 122:1--2, 3--4, 4--5, 6--7, 8--9.
Romans 13:11--14. Matthew 24:37--44.
\item[Year B (2)] Isaiah 63:16b--17, 19b and 64:2--7.
Psalm 80:2--3, 15--16, 18--19. 1 Corinthians 1:3--9. Mark 13:33--37.
If using the Douay, look for the plea to rend the heavens at Isaiah
64:1 and the potter and clay at 64:8.
\item[Year C (3)] Jeremiah 33:14--16. Psalm 25:4--5, 8--9, 10, 14.
1 Thessalonians 3:12--4:2. Luke 21:25--28, 34--36.
\end{description}

In A two people are grinding at the same mill. Nothing in their common
work tells them when the Son of Man will come. Paul gives a way to
prepare for him before either leaves the mill. Put away jealousy and
quarreling. Take the first step toward peace with someone beside whom
you have to work. Isaiah foresees whole peoples learning to do this.
They beat their swords into implements with which to feed one another.

In B the householder going abroad gives each servant work to do. Anyone
who has sat beside a sickbed knows the difference between keeping watch
and merely staying awake. Think of someone who has entrusted something
to your care. Paul tells the Corinthians why they can persevere in such
a charge. The God who called them into fellowship with his Son is faithful.

In C the disciples hear Jesus name things that make people faint with
fear. He tells them when to stand up and lift their heads. They can do
it because their redemption is near. Paul asks that their love for one
another grow while they wait. Pray for the person whose voice you would
most want to hear in a frightening hour. If that person is alive, let
them hear yours today.

St Cyril of Jerusalem once prepared his hearers to profess that Christ
will come again. He sets the child wrapped in swaddling clothes beside
the Judge clothed in light. At each coming Christians can go out to meet
the same Lord (\href{https://www.newadvent.org/fathers/310115.htm}{\emph{Catechetical Lecture} 15.1}).
Begin the week with the psalmist's petition. “Shew, O Lord, thy ways
to me, and teach me thy paths” (Douay Psalm 24:4b, U.S. Psalm 25:4).
Say it at the beginning of the day's work.

\subsection{Second Sunday. A voice beside the Jordan}

\begin{description}
\item[Year A (4)] Isaiah 11:1--10. Psalm 72:1--2, 7--8, 12--13, 17.
Romans 15:4--9. Matthew 3:1--12.
\item[Year B (5)] Isaiah 40:1--5, 9--11. Psalm 85:9--10, 11--12, 13--14.
2 Peter 3:8--14. Mark 1:1--8.
\item[Year C (6)] Baruch 5:1--9. Psalm 126:1--2, 2--3, 4--5, 6.
Philippians 1:4--6, 8--11. Luke 3:1--6.
\end{description}

In A John asks the people coming for baptism to show the fruit of
repentance. Isaiah describes the rule of the promised son of Jesse
in terms a poor man could recognize. His case will get a fair hearing.
Paul asks believers to welcome one another as Christ has welcomed
them. Give someone whose side you have never heard time to speak.

In B Isaiah's herald has news that a frightened city can bear to hear.
Its Lord comes in strength, and those strong arms gather lambs.
John directs his hearers to the stronger one who will baptize them in
the Holy Spirit. Peter, writing to people who find the Lord slow,
discerns his patience with sinners. Take the time he is giving you to
name a wrong and seek forgiveness.

In C Luke names an emperor, a governor, rulers and high priests. Then
the word of God comes to John in the desert. Hear Baruch's Jerusalem
called to climb to the height and look east. The children she had seen
led away on foot are returning. In prayer name one household still
waiting for someone to come home.

In her \href{https://ccel.org/ccel/teresa/life/life.viii.x.html}{\emph{Life} 9.4--5},
St Teresa of Ávila remembers wanting to keep Jesus company in the garden
when he was alone. She knew how hard it was to keep her thoughts there.
Begin as simply when the Baptist asks for repentance. Stay a moment with
Christ before you begin to explain yourself. Pray, “Shew us, O Lord,
thy mercy; and grant us thy salvation” (Douay Psalm 84:8, U.S. Psalm
85:8). Use the plural when you pray for someone you find hard to forgive.

\subsection{Third Sunday. What shall we do?}

\begin{description}
\item[Year A (7)] Isaiah 35:1--6a, 10. Psalm 146:6--7, 8--9, 9--10.
James 5:7--10. Matthew 11:2--11. The response to the psalm may also
be Alleluia.
\item[Year B (8)] Isaiah 61:1--2a, 10--11. Luke 1:46--48, 49--50, 53--54
(Mary's canticle). 1 Thessalonians 5:16--24. John 1:6--8, 19--28.
\item[Year C (9)] Zephaniah 3:14--18a. Isaiah 12:2--3, 4, 5--6
(a canticle). Philippians 4:4--7. Luke 3:10--18.
\end{description}

In A John's question comes from prison. Jesus gives his messengers news
they can carry past the prison door. People see and walk, the dead are
raised, and poor people hear the gospel. St John Chrysostom reads the
episode from the messengers' side. John, he thinks, wants his own
disciples to come to Jesus and see these things for themselves
(\href{https://www.newadvent.org/fathers/200136.htm}{\emph{Homily on Matthew} 36.2--3}).
Send someone confined at home a piece of good news you have seen.

In B the men questioning John want a name for him. He calls himself a
voice and speaks of the one already standing among them. Mary's song
answers Isaiah's good news to the poor with thanksgiving for the hungry
whom God has fed. Bring food to someone who needs it. In prayer remember
the person who first made sure that you were fed.

In C the crowds, tax collectors and soldiers each ask John what they
should do. His answers reach into the linen closet and the paycheck. Justice begins with a coat someone else can wear and with the
money a person can take home free from extortion. Look over one
transaction in which another person's livelihood rests in your hands.

Paul can ask for thanksgiving in the very act of asking for help.
Chrysostom observes that a penitent can grieve over sin and rejoice in
Christ who forgives it (\href{https://www.newadvent.org/fathers/230214.htm}{\emph{Homily on Philippians} 14, on 4:4--7}).
For Gaudete Sunday pray the opening of Mary's song. “My soul doth
magnify the Lord. And my spirit hath rejoiced in God my Saviour”
(Douay Luke 1:46--47). Bring one unfinished sorrow with you as you
thank him for a kindness.

\subsection{Fourth Sunday. Receiving him}

\begin{description}
\item[Year A (10)] Isaiah 7:10--14. Psalm 24:1--2, 3--4, 5--6.
Romans 1:1--7. Matthew 1:18--24.
\item[Year B (11)] 2 Samuel 7:1--5, 8b--12, 14a, 16.
Psalm 89:2--3, 4--5, 27, 29. Romans 16:25--27. Luke 1:26--38.
\item[Year C (12)] Micah 5:1--4a. Psalm 80:2--3, 15--16, 18--19.
Hebrews 10:5--10. Luke 1:39--45. The Bethlehem prophecy is Micheas
5:2--5a in the Douay.
\end{description}

In A Joseph has already decided how to spare Mary public disgrace when
the angel speaks to him. Matthew lets us see him wake and take Mary
into his home. The child she carries is the promised Emmanuel.
Consider what it would mean to give someone shelter under your good
name. Keep one confidence that would cost another person dearly if
repeated.

In B David has a house of cedar and wants to build the Lord a house.
God reminds him who first brought him in from the pasture, then
promises a house that will endure. Gabriel brings that promise to a
woman in Nazareth. St Bernard of Clairvaux invites his hearer to stand
in the silence before Mary's answer. The angel waits to return to God,
and those who need a Saviour wait with him
(\href{https://archive.org/details/sermonsofstberna00bernuoft/page/68/mode/1up}{\emph{Homilies on the Missus est} 4.8}).
Allow a moment for her words before going on to yours.

In C Elizabeth is the first person Luke lets us hear call Mary the mother
of her Lord. Even the child Elizabeth bears shares in the greeting.
Beside this meeting the Church reads from Hebrews about the body the
Son receives and offers in obedience. Carry your own prayer to someone's
door. A visit that frees a tired parent to rest can be a way of receiving
the child for whom you are preparing.

Pray Mary's answer, “Behold the handmaid of the Lord: be it done to me
according to thy word” (Douay Luke 1:38). Return to it before the next
ordinary task for which someone is depending on you. When this Sunday
is December 24, the Advent Sunday has its own Mass. At the evening
celebration of the Nativity, turn to \hyperref[sec:pc-christmas]{Christmas}.
